\begin{lemma}
\label{lemma-finite-length}
Let $R$ be a domain with fraction field $K$.
Let $M$ be an $R$-submodule of $K^{\oplus r}$.
Assume $R$ is local Noetherian of dimension $1$.
For any nonzero $x \in R$ we have $\text{length}_R(R/xR) < \infty$
and
$$
\text{length}_R(M/xM) \leq r \cdot \text{length}_R(R/xR).
$$
\end{lemma}

\begin{proof}
If $x$ is a unit then the result is true. Hence we may assume
$x \in \mathfrak m$ the maximal ideal of $R$. Since $x$ is not
zero and $R$ is a domain we have $\dim(R/xR) = 0$, and hence
$R/xR$ has finite length. Consider $M \subset K^{\oplus r}$ as
in the lemma. We may assume that the elements of $M$ generate
$K^{\oplus r}$ as a $K$-vector space after replacing $K^{\oplus r}$
by a smaller subspace if necessary.

\medskip\noindent
Suppose first that $M$ is a finite $R$-module. In that case we can clear
denominators and assume $M \subset R^{\oplus r}$. Since
$M$ generates $K^{\oplus r}$ as a vectors space we see that
$R^{\oplus r}/M$ has finite length. In particular there exists
an integer $c \geq 0$ such that $x^cR^{\oplus r} \subset M$.
Note that $M \supset xM \supset x^2M \supset \ldots$ is a sequence of
modules with successive quotients each isomorphic to $M/xM$. Hence
we see that
$$
n \text{length}_R(M/xM) = \text{length}_R(M/x^nM).
$$
The same argument for $M = R^{\oplus r}$ shows that
$$
n \text{length}_R(R^{\oplus r}/xR^{\oplus r}) =
\text{length}_R(R^{\oplus r}/x^nR^{\oplus r}).
$$
By our choice of $c$ above we see that $x^nM$ is sandwiched between
$x^n R^{\oplus r}$ and $x^{n + c}R^{\oplus r}$. This easily gives that
$$
r(n + c) \text{length}_R(R/xR)
\geq
n \text{length}_R(M/xM)
\geq
r (n - c) \text{length}_R(R/xR)
$$
Hence in the finite case we actually get the result of the lemma with
equality.

\medskip\noindent
Suppose now that $M$ is not finite. Suppose that the length of $M/xM$ is
$\geq k$ for some natural number $k$. Then we can find
$$
0 \subset N_0 \subset N_1 \subset N_2 \subset \ldots N_k \subset M/xM
$$
with $N_i \not = N_{i + 1}$ for $i = 0, \ldots k - 1$.
Choose an element $m_i \in M$ whose congruence class mod $xM$ falls
into $N_i$ but not into $N_{i - 1}$ for $i = 1, \ldots, k$.
Consider the finite $R$-module $M' = Rm_1 + \ldots + Rm_k \subset M$.
Let $N'_i \subset M'/xM'$ be the inverse image of $N_i$.
It is clear that $N'_i \not =N'_{i + 1}$ by our choice of $m_i$.
Hence we see that $\text{length}_R(M'/xM') \geq k$. By the
finite case we conclude $k \leq r\text{length}_R(R/xR)$
as desired.
\end{proof}
